\section{Verification purpose}

This case verifies spatial moisture lookup, moisture damping, and extinction behavior in otherwise uniform fuel.  The verification question is whether the four regional spread rates meet their moisture-specific references, including negligible spread at the extinction condition.

The scientific setup and historical expected response are reconstructed from Section~3.2 of the supplied \emph{ELMFIRE Guide}. The Guide is used as reference provenance, not as an instruction or a substitute for evidence. Configuration statements below are checked against the generated GeoTIFFs and namelist; current results and decisions come only from case-specific outputs and generated metrics. Where the generated experiment differs from the historical example, the report states that difference explicitly.

\section{Mathematical formulation and numerical implementation}

Fuel moisture reduces reaction intensity through the model's moisture-damping term; at or above the fuel's moisture of extinction, sustainable spread approaches zero.  Only the one-hour moisture raster varies, so regional velocity differences should follow that damping response.

For a dead-fuel moisture ratio \(r=M_f/M_x\), the Rothermel damping polynomial used in the configured range is
\begin{equation}
\eta_M(r)=1-2.59r+5.11r^2-3.52r^3, \qquad 0\le r\le1.
\end{equation}
Reaction intensity contains this damping factor, \(I_R=\Gamma^{\prime}(w_{n,d}\eta_{M,d}+w_{n,l}\eta_{M,l})\eta_s h\), and the no-wind rate is \(R_0=I_R\xi/(\rho_b\epsilon Q_{ig})\). At or above extinction, the dead-fuel contribution cannot sustain propagation.

\section{Derivation of expected behavior}

For Fuel Model 3 the configured moisture values 6, 12, 18, and 25 percent correspond to reference rates \(1.511\), \(1.089\), \(0.801\), and \(0~\mathrm{m\,min^{-1}}\).  The rates must decrease monotonically as moisture rises.

% BEGIN EXPLICIT METRIC DERIVATION
The dead one-hour moisture ratio is varied while all other generated fuel
and moisture properties remain fixed.  Applying
\(R_0=I_R\xi/(\rho_b\epsilon Q_{ig})\) and
\(Q_{ig}=250+1116M_f\) gives
\[
\begin{array}{c|cc}
M_{1h}&I_R~(\mathrm{kW\,m^{-2}})&R_0~(\mathrm{m\,min^{-1}})\\
\hline
0.06&549.174&1.511830\\
0.12&479.466&1.089719\\
0.18&413.934&0.801066\\
0.25&\simeq3.9\times10^{-13}&\simeq6.5\times10^{-16}
\end{array}
\]
so the report references are \(1.511,\ 1.089,\ 0.801,\) and
\(0~\mathrm{m\,min^{-1}}\).  Nonzero rows use
\(e_R=|R_{\max}-R_{\mathrm{ref}}|/R_{\mathrm{ref}}\); the extinguished
row uses the absolute residual \(|R_{\max}|\), avoiding division by a
zero reference.
% END EXPLICIT METRIC DERIVATION

\section{Verification metrics and acceptance criteria}

The maximum finite spread rate in each quadrant measures the local head response without averaging across moisture boundaries. Nonzero targets use 10 percent relative tolerance. The extinction quadrant uses an absolute \(0.05~\mathrm{m\,min^{-1}}\) limit because relative error is undefined at zero. All four components must pass.

\section{Simulation configuration and predicted outcome}

The grid uses uniform Fuel Model 3, flat terrain, zero wind, central ignition, and one-hour moisture quadrants of 6, 12, 18, and 25 percent. Predicted propagation slows monotonically and stops in the 25-percent quadrant.

\begin{center}
\begin{tabular}{@{}p{0.20\textwidth}p{0.20\textwidth}p{0.12\textwidth}p{0.38\textwidth}@{}}
\toprule
Parameter & Configured value & Units & Role in the experiment \\
\midrule
Grid & \(126\times126\) at 30 & m per cell & Resolves the whole spatial response on the generated north-up raster \\
Projection & EPSG:32635 & --- & Fixes the map-coordinate and row-orientation convention \\
Time control & \(\Delta t=1\), \(t_{stop}=259200\) & s & Whole-second step; stop time is an exact timestep multiple \\
Meteorology interval & 3600 & s & Defines weather-band timing and stalled-final compatibility \\
Live moisture & herb 60; woody 90 & percent & Holds live-fuel damping fixed while the named mechanism is isolated \\
Evidence & prepared inputs; final case outputs & --- & Separates configuration provenance from acceptance evidence \\
\bottomrule
\end{tabular}
\end{center}

Figure~\ref{fig:input-configuration} shows the prescribed spatial condition used in this experiment. The displayed field represents the initial level-set geometry or a controlling fuel, terrain, moisture, wind, or canopy quantity, as identified by the figure labels. The plotted values come from the prepared simulation inputs, not from the expected response.

\begin{figure}[H]
  \centering
  \IfFileExists{../figures/input_configuration.pdf}{%
    \includegraphics[width=0.98\textwidth,height=0.46\textheight,keepaspectratio]{input_configuration.pdf}%
  }{\fbox{\parbox{0.88\textwidth}{Prepared-input evidence is unavailable. The prescribed configuration is described in the text; this absence does not establish that it was simulated.}}}
  \caption{Prepared input conditions from the prescribed spatial fields. Coordinates are in metres; the colorbar identifies the displayed field and its physical units or category. The ignition panel is derived directly from \(\phi\leq0\).}
  \label{fig:input-configuration}
\end{figure}

\section{Actual simulation results}

The required evidence is generated only from case-specific ELMFIRE outputs. The current postprocessing status is \textbf{\VerificationStatus}; \PassingMetricCount{} of \MetricCount{} required metric rows pass across \VariantCount{} simulation condition(s). A missing output is reported as NOT RUN or NOT EVALUABLE and is not interpreted as verification evidence.

\begin{figure}[H]
  \centering
  \IfFileExists{../figures/whole_domain_result.pdf}{%
    \includegraphics[width=0.82\textwidth,height=0.58\textheight,keepaspectratio]{whole_domain_result.pdf}%
  }{\fbox{\parbox{0.85\textwidth}{No case-specific whole-domain ELMFIRE raster is available.}}}
  \caption{Final whole-domain ELMFIRE time of arrival for the displayed simulation condition, in seconds. The coordinate reference system is stated in the figure. The calculated metrics, not visual appearance, determine acceptance.}
\end{figure}

\section{Calculated metrics and verification decision}

\begin{center}
\begin{tabular}{@{}llll@{}}
\toprule
Simulation condition / metric & Expected & Calculated & Status \\
\midrule
\MetricRows
\bottomrule
\end{tabular}
\end{center}

The displayed overall decision is \textbf{\VerificationStatus}. When calculated verification and validation metrics exists, result analysis procedure generates every table value and status; when it is absent, the report states NOT EVALUABLE with no invented rows.

\section{Technical interpretation}

Nonmonotone behavior suggests moisture raster or damping errors; significant spread in the extinction quadrant is a direct failure. No scientific conclusion should be drawn when the status is NOT RUN or NOT EVALUABLE.

\section{Scope and limitations}

Only one fuel model and the one-hour dead-fuel moisture field are varied. This case provides implementation verification evidence only.
